\documentclass[11pt,a4paper]{article}
\usepackage[UTF8]{ctex}
\usepackage[margin=2.5cm]{geometry}
\usepackage{booktabs}
\usepackage{longtable}
\usepackage{listings}
\usepackage{xcolor}
\usepackage{hyperref}
\usepackage{enumitem}
\usepackage{amssymb}
\hypersetup{colorlinks=true,linkcolor=blue,urlcolor=blue}

\lstset{
  basicstyle=\ttfamily\small,
  breaklines=true,
  frame=single,
  backgroundcolor=\color{gray!10},
  keywordstyle=\color{blue},
  commentstyle=\color{gray},
  language=bash
}

\title{\textbf{LLM Context Compaction Proxy}\\[0.3em]
\large Never Full, Never Stop — 永不超限,永不停歇\\[0.3em]
\large README(LaTeX 版)}
\author{LLM Context Compaction Proxy Contributors}
\date{2026 年 8 月}

\begin{document}

\maketitle

% ============================================================
\section{简介}
% ============================================================

LLM Context Compaction Proxy 是一个透明、零配置的 LLM API 上下文压缩代理。
将其置于 AI Agent 与任意 LLM 提供商之间,当对话上下文逼近模型 token 上限时,
它会自动压缩对话历史,使 Agent 无限期运行而不触碰上下文窗口墙。

\emph{40 项压缩技术}横跨 7 代演进,来源涵盖学术研究
(ARC、AFM、PACMS、CoMem、MemSkill)与生产实践
(Cursor、Devin、SWE-agent)。

% ============================================================
\section{Feature Highlights}
% ============================================================

\subsection{Context Compression Engine}

\begin{longtable}{@{}lp{9.5cm}@{}}
  \toprule
  Feature & Description \\
  \midrule
  Preemptive Compaction & Triggers at 80\% context before hitting the limit \\
  Overflow Recovery & Compress + retry on context overflow errors \\
  Incremental Compaction & Only summarize new messages since last compaction (CoMem) \\
  Parallel Compaction & Split large contexts into parallel blocks \\
  Parallel Summary Merge & Merge parallel block summaries into one cohesive summary \\
  Dual-Layer Compression & Gateway (85\% reduction) + Agent (50\% of L1), three fallback paths \\
  CJK-Aware Token Estimation & Accurate CJK/ASCII/other token counting (two-stage) \\
  Compaction Safety Verification & Compacted result must be smaller than original \\
  Smart Truncation & Role-based budget allocation (head/tail split) \\
  Identifier Preservation & Keep code identifiers intact during compression \\
  Mandatory Identifier List & Force-preserve UUIDs/hashes/URLs/file names verbatim \\
  Compaction Cache & 30-min TTL, instruction-aware salt keys \\
  Compaction-Item Pruning & Prune low-value artifacts left by previous compactions \\
  Query Placement Optimization & Query-at-end, tool-pair adjacency, recent-window reorder \\
  Cache-Optimized Ordering & Layer-hash-aware reorder + \texttt{cache\_control} breakpoints \\
  Structure-Aware Tool Output Compression & Dedicated compressors for code / logs / JSON \\
  Manual Compaction (\texttt{/compact}) & Force compact any session, like Claude Code's /compact \\
  Aggressive Truncation Fallback & Degrade to truncation when compaction fails or grows \\
  \bottomrule
\end{longtable}

\subsection{Memory \& Session}

\begin{longtable}{@{}lp{9.5cm}@{}}
  \toprule
  Feature & Description \\
  \midrule
  Semantic Memory & Episodic-semantic dual-layer memory, 6 knowledge slots \\
  Cross-Session Memory + FTS5 & SQLite session persistence with full-text search \\
  User Profile Memory & Persistent user preferences (USER.md), bounded persistence \\
  Session Transcript Archive & Save/restore raw transcripts for session resume \\
  Session Resume & Resume from transcript + summary (\texttt{POST /sessions/\{id\}/resume}) \\
  Background Knowledge Injection & Inject recent session knowledge into compaction prompts \\
  Session Context Injection & Auto-inject session context into outgoing requests \\
  ARC References & Replace lengthy tool results with ID references \\
  ARC Persistent Recall & Persist ARC entries to DB, recall via \texttt{GET /arc/\{arc\_id\}} \\
  Commitment Extraction (CCL) & Extract and preserve goals/constraints/decisions/errors \\
  LLM Memory Extraction & LLM-driven knowledge extraction with regex fallback \\
  Thought Masking & Strip reasoning\_content to save tokens \\
  Secret Redaction & Auto-redact API keys/JWT/passwords/tokens \\
  Orphan Tool Pair Sanitization & Clean up dangling tool\_call/tool\_result pairs \\
  Tool Pair Adjacency Fix & Repair non-adjacent tool call/result ordering \\
  \bottomrule
\end{longtable}

\subsection{Provider \& Protocol}

\begin{longtable}{@{}lp{9.5cm}@{}}
  \toprule
  Feature & Description \\
  \midrule
  Provider Abstraction Layer & Universal OpenAI/Anthropic/Gemini compatibility (ABC) \\
  Auto Provider Detection & Detect format from model name / headers / URL \\
  OpenAI$\rightarrow$Anthropic Request Conversion & Full field mapping incl. tool\_calls $\rightarrow$ tool\_use \\
  Anthropic$\rightarrow$OpenAI Response Conversion & Content/thinking/tool\_use mapped back to OpenAI \\
  SSE Streaming Conversion & Anthropic SSE $\rightarrow$ OpenAI SSE event-by-event \\
  Anthropic Thinking-Only Retry & Auto streaming retry when response has only empty thinking \\
  Separate Compaction Provider & Independent upstream \& API key for compaction model \\
  Gemini Header Auth & \texttt{x-goog-api-key} header, never in URL/logs \\
  Model Registry & 67+ models with known context limits \\
  Passthrough Routing & Any unmatched path transparently forwarded upstream \\
  \bottomrule
\end{longtable}

\subsection{Reliability \& Security}

\begin{longtable}{@{}lp{9.5cm}@{}}
  \toprule
  Feature & Description \\
  \midrule
  Circuit Breaker & 3 failures $\rightarrow$ 60s cooldown, three-state machine \\
  Thrashing Detection & Detect compaction loops (3 compacts / 5 msgs) \\
  Auth (require\_auth) & 22+ endpoints protected; loopback-only when no secret \\
  Concurrency Safety & Thread-locked semantic memory, race-free prompts \\
  Pre/Post Compaction Hooks & HTTP webhooks before/after compaction (blockable) \\
  Health \& Metrics Endpoints & \texttt{/health}, \texttt{/metrics} (Prometheus-style) \\
  Adaptive Keep-Turns & Auto-reduce kept turns when conversation is short \\
  Retry with Shrinking Window & Retry compaction with fewer turns per attempt \\
  \bottomrule
\end{longtable}

\subsection{MemSkill (Self-Evolving, V7)}

\begin{longtable}{@{}lp{9.5cm}@{}}
  \toprule
  Feature & Description \\
  \midrule
  Skill Registry & CRUD + activation lifecycle + snapshot rollback (SQLite) \\
  Skill Controller & Keyword matching + Gumbel-Top-K selection \\
  Parameter Override Pipeline & Importance weights + fidelity multipliers per skill \\
  DELETE-type Skill Pipeline Skip & Skip semantic extraction for DELETE operations \\
  Skill Performance Tracking & Reward/success stats (\texttt{/skills/\{id\}/performance}) \\
  Skill Trajectories & Recent compaction trajectories (\texttt{/skills/trajectories}) \\
  Skill Designer & Auto-design new skills (\texttt{/skills/designer/trigger}) \\
  \bottomrule
\end{longtable}

\subsection{HTTP API Endpoints (31)}

\begin{longtable}{@{}llp{7cm}@{}}
  \toprule
  Method & Endpoint & Description \\
  \midrule
  POST & \texttt{/chat/completions} (+\texttt{/v1/...}) & OpenAI-compatible chat (stream \& non-stream) \\
  POST & \texttt{/v1/messages} & Anthropic Messages API \\
  POST & \texttt{/compact} & Manual compaction of a session \\
  POST & \texttt{/summarize} & Summary-only endpoint for CompactionProvider plugins \\
  POST & \texttt{/session} & Create a session \\
  GET & \texttt{/sessions/search} & FTS5 full-text session search \\
  GET & \texttt{/sessions/recent} & Recent sessions \\
  GET & \texttt{/sessions/\{id\}} & Session detail \\
  GET & \texttt{/sessions/\{id\}/transcript} & Raw transcript export \\
  POST & \texttt{/sessions/\{id\}/resume} & Resume a session \\
  GET/POST & \texttt{/profile} & User profile get/set \\
  GET/DELETE & \texttt{/memory} & Semantic memory get/clear \\
  GET & \texttt{/arc/\{arc\_id\}} & ARC reference recall \\
  GET & \texttt{/skills}, \texttt{/skills/\{id\}} & Skill list/detail \\
  POST & \texttt{/skills} & Create skill (draft) \\
  POST & \texttt{/skills/\{id\}/activate}, \texttt{/deprecate} & Skill lifecycle \\
  POST & \texttt{/skills/\{id\}/rollback} & Rollback to snapshot version \\
  GET & \texttt{/skills/\{id\}/performance} & Skill reward stats \\
  GET & \texttt{/skills/trajectories} & Compaction trajectories \\
  POST & \texttt{/skills/designer/trigger} & Trigger skill designer \\
  GET & \texttt{/health}, \texttt{/metrics} & Health \& metrics \\
  ANY & \texttt{/\{path:.*\}} & Passthrough to upstream \\
  \bottomrule
\end{longtable}

% ============================================================
\section{工作原理}
% ============================================================

\begin{verbatim}
+-------------+     +----------------------+     +--------------+
|  AI Agent   |---->|  Compaction Proxy    |---->|  LLM Provider|
| (Claude/GPT)|<----|  localhost:8198      |<----|  (any API)   |
+-------------+     |                      |     +--------------+
                    |  - Monitor context % |
                    |  - Auto-compress     |
                    |  - Preserve key info |
                    |  - Route to provider |
                    +----------------------+
\end{verbatim}

\begin{enumerate}[leftmargin=2em]
  \item \textbf{Passthrough(透传)}——上下文低于限制的请求直接通过,零开销;
  \item \textbf{Detect(检测)}——上下文达到阈值(默认 80\%)时触发压缩;
  \item \textbf{Compress(压缩)}——使用独立压缩模型对旧消息生成摘要;
  \item \textbf{Resume(恢复)}——压缩结果替换旧消息,Agent 无缝继续。
\end{enumerate}

% ============================================================
\section{快速开始}
% ============================================================

\subsection{前置条件}

\begin{itemize}
  \item Python 3.10+
  \item \texttt{aiohttp}(\texttt{pip install aiohttp})
\end{itemize}

\subsection{安装与运行}

\begin{lstlisting}
# Clone
git clone https://github.com/your-org/llm-compaction-proxy.git
cd llm-compaction-proxy

# Set your upstream LLM API
export COMPACTION_PROXY_UPSTREAM=https://api.openai.com/v1
export COMPACTION_PROXY_MODEL=gpt-4o-mini

# Run
python3 compaction-proxy.py
\end{lstlisting}

代理监听在 \texttt{localhost:8198}。将 Agent 的 API 端点指向此处即可。

\subsection{使用环境变量文件}

\begin{lstlisting}
cp .env.example .env
# Edit .env with your settings
source .env
python3 compaction-proxy.py
\end{lstlisting}

\subsection{作为 systemd 服务}

\begin{lstlisting}
cp systemd/compaction-proxy.service ~/.config/systemd/user/
# Edit the service file to set your paths and environment
systemctl --user daemon-reload
systemctl --user enable --now compaction-proxy.service
\end{lstlisting}

% ============================================================
\section{连接 AI Agent}
% ============================================================

\subsection{Claude Code (Anthropic)}

\begin{lstlisting}
# Set the API base to the proxy
export ANTHROPIC_BASE_URL=http://localhost:8198
claude
\end{lstlisting}

或在 Claude Code 设置中:

\begin{lstlisting}[language=bash]
{
  "apiBaseUrl": "http://localhost:8198"
}
\end{lstlisting}

\subsection{OpenAI 兼容 Agent(GPT、DeepSeek、Ollama、vLLM 等)}

\begin{lstlisting}
export OPENAI_BASE_URL=http://localhost:8198/v1
export OPENAI_API_KEY=your-key
\end{lstlisting}

适用于任何使用 OpenAI SDK 或兼容 API 的 Agent。

\subsection{Cursor / Continue / 其他 IDE 扩展}

在扩展设置中设置 API base URL:
\begin{itemize}
  \item \textbf{Cursor}:Settings $\rightarrow$ Models $\rightarrow$ OpenAI API Base
        $\rightarrow$ \texttt{http://localhost:8198/v1}
  \item \textbf{Continue}:\texttt{config.json} $\rightarrow$
        \texttt{"apiBase": "http://localhost:8198/v1"}
\end{itemize}

\subsection{自定义 Agent(直接 API 调用)}

\begin{lstlisting}[language=Python]
import openai

client = openai.OpenAI(
    base_url="http://localhost:8198/v1",
    api_key="your-key"
)

response = client.chat.completions.create(
    model="gpt-4o-mini",
    messages=[{"role": "user", "content": "Hello!"}],
    stream=True
)
\end{lstlisting}

% ============================================================
\section{支持的提供商}
% ============================================================

\begin{longtable}{@{}llc@{}}
  \toprule
  提供商 & 格式 & 自动检测 \\
  \midrule
  OpenAI & OpenAI & \checkmark \\
  Anthropic & Anthropic & \checkmark \\
  Google Gemini & Gemini & \checkmark \\
  DeepSeek & OpenAI & \checkmark \\
  Ollama & OpenAI & \checkmark \\
  vLLM & OpenAI & \checkmark \\
  LiteLLM & OpenAI & \checkmark \\
  OpenRouter & OpenAI & \checkmark \\
  Together AI & OpenAI & \checkmark \\
  Groq & OpenAI & \checkmark \\
  Fireworks & OpenAI & \checkmark \\
  MaaS 代理 & Anthropic & \checkmark \\
  \bottomrule
\end{longtable}

设置 \texttt{COMPACTION\_PROXY\_UPSTREAM\_IS\_ANTHROPIC=auto}(默认)自动检测,
或使用 \texttt{true}/\texttt{false} 强制覆盖。

% ============================================================
\section{接入你的 LLM 提供商}
% ============================================================

代理位于 Agent 与任意 LLM API 之间。\textbf{一个变量决定上游}:
\texttt{COMPACTION\_PROXY\_UPSTREAM}。其余均为可选。

\subsection{OpenAI / OpenAI 兼容系}

\begin{lstlisting}
export COMPACTION_PROXY_UPSTREAM=https://api.openai.com/v1
export COMPACTION_PROXY_MODEL=gpt-4o-mini
python3 compaction-proxy.py
\end{lstlisting}

开箱即用——代理转发请求,并将压缩结果转换回 Agent 期望的 OpenAI 格式。

\subsection{Anthropic (Claude)}

\begin{lstlisting}
export COMPACTION_PROXY_UPSTREAM=https://api.anthropic.com
export COMPACTION_PROXY_UPSTREAM_IS_ANTHROPIC=true
export COMPACTION_PROXY_MODEL=claude-sonnet-4-5
python3 compaction-proxy.py
\end{lstlisting}

代理使用 Anthropic Messages API(\texttt{/v1/messages}),Anthropic 原生 Agent
可直接连接;OpenAI 格式 Agent 指向同一代理时自动应用格式转换。

\subsection{Google Gemini}

\begin{lstlisting}
export COMPACTION_PROXY_UPSTREAM=https://generativelanguage.googleapis.com
export COMPACTION_PROXY_MODEL=gemini-2.0-flash
python3 compaction-proxy.py
\end{lstlisting}

API key 通过 \texttt{x-goog-api-key} 请求头传递(绝不进入 URL),不会出现在访问日志中。

\subsection{DeepSeek}

\begin{lstlisting}
export COMPACTION_PROXY_UPSTREAM=https://api.deepseek.com/v1
export COMPACTION_PROXY_MODEL=deepseek-chat
python3 compaction-proxy.py
\end{lstlisting}

DeepSeek 上下文窗口(\texttt{deepseek-chat}/\texttt{coder}/\texttt{r1}/\texttt{v3},
128K)已预注册,无需额外配置即可完成压力预估。

\subsection{本地模型 (Ollama)}

\begin{lstlisting}
export COMPACTION_PROXY_UPSTREAM=http://localhost:11434/v1   # default
export COMPACTION_PROXY_MODEL=llama3.1
python3 compaction-proxy.py
\end{lstlisting}

\subsection{验证是否生效}

\begin{lstlisting}
# Health check
curl http://localhost:8198/health

# Send a test request through the proxy (OpenAI format)
curl -X POST http://localhost:8198/v1/chat/completions \
  -H "Content-Type: application/json" \
  -H "Authorization: Bearer YOUR_KEY" \
  -d '{"model":"gpt-4o-mini","messages":[{"role":"user","content":"Hello"}]}'
\end{lstlisting}

\subsection{使用提示}

\begin{itemize}
  \item \textbf{压缩模型 $\neq$ 主模型}:设置
        \texttt{COMPACTION\_PROXY\_COMPACTION\_UPSTREAM} /
        \texttt{COMPACTION\_PROXY\_COMPACTION\_API\_KEY},可在主模型之外
        使用更廉价的模型做摘要;
  \item \textbf{Anthropic 格式 MaaS 代理}(如讯飞 coding API):设置
        \texttt{COMPACTION\_PROXY\_UPSTREAM\_IS\_ANTHROPIC=true} 以应用请求转换;
  \item \textbf{切勿在配置中硬编码密钥}:使用环境变量或 \texttt{apiKeyEnv} 模式,
        代理从不记录 API key。
\end{itemize}

% ============================================================
\section{40 项压缩技术}
% ============================================================

\subsection{V3 核心(8 项)}

\begin{longtable}{@{}clp{8cm}@{}}
  \toprule
  \# & 技术 & 说明 \\
  \midrule
  1  & \textbf{ARC} 地址化引用 & 用紧凑的 \texttt{@ref} 指针替换重复内容 \\
  2  & \textbf{AFM} 自适应保真 & 消息级 Full/Compressed/Placeholder + 子模增强 \\
  3  & \textbf{PACMS} 子模选择 & 贪心、保真感知的消息选择 \\
  4  & \textbf{CCL} 承诺提取 & 提取并保留对话中的承诺 \\
  5  & 抖动检测 & 检测压缩循环(3 次压缩 / 5 条消息) \\
  6  & 两阶段 Token 估算 & CJK 感知的精确 token 计数 \\
  7  & 预飞安全验证 & 发送前验证上下文合法性 \\
  8  & 并行压缩 & 大上下文分块并行压缩 \\
  \bottomrule
\end{longtable}

\subsection{V4 增强(10 项)}

\begin{longtable}{@{}clp{8cm}@{}}
  \toprule
  \# & 技术 & 来源 \\
  \midrule
  9  & 情节—语义双层记忆 & arXiv:2605.17625 \\
  10 & 思考掩码 & Kevin-32B / Devin 模式 \\
  11 & 标签选择性保留 & SWE-agent / memor-ai \\
  12 & 结构感知工具输出压缩 & memor-ai \\
  13 & 密钥脱敏 & Cursor / memor-ai \\
  14 & 缓存优化消息排序 & 层哈希感知 + \texttt{cache\_control} 断点 \\
  15 & 增量压缩 & CoMem (arXiv:2605.30842) \\
  16 & 四信号记忆评分 & 语义 50\%、近因 25\%、类型 15\%、质量 10\% \\
  17 & 查询位置优化 & 查询置尾、配对邻接、近期窗口 \\
  18 & 压缩项剪枝 & OpenAI 模式 \\
  \bottomrule
\end{longtable}

\subsection{V5 会话(5 项)}

\begin{longtable}{@{}clp{8cm}@{}}
  \toprule
  \# & 技术 & 说明 \\
  \midrule
  19 & 跨会话记忆 + FTS5 & SQLite 会话持久化 + 全文搜索 \\
  20 & 压缩前后 Hooks & 压缩前后的 HTTP webhook \\
  21 & 孤儿工具对清理 & 清理悬空的 tool call/result 对 \\
  22 & 用户画像记忆 & 持久化用户偏好(USER.md) \\
  23 & 可插拔压缩引擎 & 切换压缩后端(default / dual-layer) \\
  \bottomrule
\end{longtable}

\subsection{V6 Provider(5 项)}

\begin{longtable}{@{}clp{8cm}@{}}
  \toprule
  \# & 技术 & 说明 \\
  \midrule
  24 & Provider 抽象层 & 通用 OpenAI/Anthropic/Gemini 兼容 \\
  25 & 压缩 Provider & 独立的上游与 API key \\
  26 & Provider 专属溢出检测 & 各 Provider 错误模式匹配 \\
  27 & 双层压缩 & 网关(85\% 缩减)+ Agent(L1 的 50\%) \\
  28 & CachedSystemPrompt & 10 层前缀缓存 + \texttt{cache\_control} 断点 \\
  \bottomrule
\end{longtable}

\subsection{V7 MemSkill(5 项)}

\begin{longtable}{@{}clp{8cm}@{}}
  \toprule
  \# & 技术 & 来源 \\
  \midrule
  36 & 自演进记忆技能 & arXiv:2602.02474 \\
  37 & 技能控制器 & 关键词匹配 + Gumbel-Top-K 选择 \\
  38 & 技能注册表 & CRUD + 激活生命周期 + 快照回滚 \\
  39 & 参数覆盖管线 & 重要性权重 + 保真乘数 \\
  40 & DELETE 型技能管线跳过 & 跳过 DELETE 操作的语义提取 \\
  \bottomrule
\end{longtable}

\subsection{基础设施}

\begin{longtable}{@{}lp{8cm}@{}}
  \toprule
  特性 & 说明 \\
  \midrule
  自动 Provider 检测 & 自动识别 OpenAI vs Anthropic vs Gemini \\
  模型注册表 & 67+ 已知上下文窗口的模型 \\
  熔断器 & 3 次失败 $\rightarrow$ 60 秒冷却 \\
  压缩缓存 & 30 分钟 TTL,避免对未变化上下文重复压缩 \\
  预压缩 & 达到限制前 80\% 触发 \\
  溢出恢复 & 溢出错误时压缩并重试 \\
  智能截断 & 基于角色的预算分配 \\
  标识符保留 & 压缩时保持代码标识符完整 \\
  摘要合并 & 合并并行压缩块结果 \\
  \bottomrule
\end{longtable}

% ============================================================
\section{配置说明}
% ============================================================

所有配置均为带合理默认值的环境变量,完整列表见 \texttt{.env.example}。

\subsection{核心配置}

\begin{longtable}{@{}llp{7cm}@{}}
  \toprule
  变量 & 默认值 & 说明 \\
  \midrule
  \texttt{COMPACTION\_PROXY\_PORT} & \texttt{8198} & 代理监听端口 \\
  \texttt{COMPACTION\_PROXY\_UPSTREAM} & \texttt{http://localhost:11434/v1} & 上游 LLM API 基地址 \\
  \texttt{COMPACTION\_PROXY\_UPSTREAM\_IS\_ANTHROPIC} & \texttt{auto} & 强制 Anthropic 格式 \\
  \texttt{COMPACTION\_PROXY\_MODEL} & \texttt{gpt-4o-mini} & 压缩/摘要模型 \\
  \texttt{COMPACTION\_PROXY\_DEFAULT\_MODEL} & \texttt{gpt-4o-mini} & 无模型请求的默认模型 \\
  \bottomrule
\end{longtable}

\subsection{压缩行为}

\begin{longtable}{@{}llp{7cm}@{}}
  \toprule
  变量 & 默认值 & 说明 \\
  \midrule
  \texttt{COMPACTION\_PROXY\_KEEP\_TURNS} & \texttt{6} & 保留的最近轮次(不压缩) \\
  \texttt{COMPACTION\_PROXY\_MAX\_RETRIES} & \texttt{3} & 压缩 API 最大重试次数 \\
  \texttt{COMPACTION\_PROXY\_TIMEOUT} & \texttt{120} & 压缩超时(秒) \\
  \texttt{COMPACTION\_PROXY\_PARALLEL\_BLOCKS} & \texttt{3} & 并行压缩块数 \\
  \texttt{COMPACTION\_PROXY\_PREEMPTIVE\_THRESHOLD} & \texttt{0.80} & 上下文 80\% 时预压缩 \\
  \bottomrule
\end{longtable}

\subsection{双层压缩}

\begin{longtable}{@{}llp{7cm}@{}}
  \toprule
  变量 & 默认值 & 说明 \\
  \midrule
  \texttt{COMPACTION\_PROXY\_ENGINE} & \texttt{default} & 压缩引擎(default / dual-layer) \\
  \texttt{COMPACTION\_PROXY\_DUAL\_GATEWAY\_RATIO} & \texttt{0.15} & L1 网关:保留 15\%(85\% 缩减) \\
  \texttt{COMPACTION\_PROXY\_DUAL\_AGENT\_RATIO} & \texttt{0.50} & L2 Agent:保留 L1 输出的 50\% \\
  \bottomrule
\end{longtable}

\subsection{安全}

\begin{longtable}{@{}llp{7cm}@{}}
  \toprule
  变量 & 默认值 & 说明 \\
  \midrule
  \texttt{COMPACTION\_PROXY\_API\_SECRET} & \textit{空} & 代理访问密钥 \\
  \texttt{COMPACTION\_PROXY\_REDACT\_SECRETS} & \texttt{1} & 自动脱敏消息中的密钥 \\
  \bottomrule
\end{longtable}

\subsection{会话与记忆}

\begin{longtable}{@{}llp{7cm}@{}}
  \toprule
  变量 & 默认值 & 说明 \\
  \midrule
  \texttt{COMPACTION\_PROXY\_BACKGROUND\_SESSIONS} & \texttt{3} & 并行压缩的后台会话数 \\
  \texttt{COMPACTION\_PROXY\_LLM\_MEMORY} & \texttt{1} & 启用 LLM 记忆提取 \\
  \texttt{COMPACTION\_PROXY\_MEMSKILL} & \texttt{0} & 启用自演进 MemSkill \\
  \bottomrule
\end{longtable}

% ============================================================
\section{架构}
% ============================================================

\begin{verbatim}
Request Flow:
  Agent -> Proxy -> [context check] -> Upstream
                         |
                    if > 80% full:
                         |
                    +----v----+
                    | Compact |--> Compaction Model (separate provider)
                    | Engine  |<-- Compressed Summary
                    +----+----+
                         |
                    Replace old messages
                    with compressed summary
                         |
                         v
                    -> Upstream (with compacted context)
\end{verbatim}

\subsection{双层压缩 (V6)}

\begin{verbatim}
Original Context (100%)
       |
  +----v----+
  |  L1 Gateway  |  Keep 15% -> 85% reduction
  |  Compression |
  +----+----+
       | 15% retained
  +----v----+
  |  L2 Agent   |  Keep 50% of L1 -> 92.5% total reduction
  |  Compression |
  +----+----+
       | 7.5% of original
       v
  Final Context
\end{verbatim}

% ============================================================
\section{API 兼容性}
% ============================================================

代理是透明的——转发所有 API 端点并保持原始请求/响应格式。
Agent 无需知道它的存在。

\subsection{OpenAI 格式}

\begin{verbatim}
POST /v1/chat/completions
POST /v1/models
\end{verbatim}

\subsection{Anthropic 格式}

\begin{verbatim}
POST /v1/messages
POST /v1/messages/count_tokens
\end{verbatim}

\subsection{Gemini 格式}

\begin{verbatim}
POST /v1beta/models/{model}:generateContent
POST /v1beta/models/{model}:streamGenerateContent
\end{verbatim}

% ============================================================
\section{许可证}
% ============================================================

MIT License——详见 \texttt{LICENSE} 文件。

\end{document}

